% Insert this text in the Data collection section, after the paragraph
% describing source preservation.

For every LLM-generated record, the final index also stores explicit provenance
metadata: \texttt{source\_site}, \texttt{source\_url},
\texttt{original\_title}, and \texttt{collection\_date}. These fields are
populated deterministically from the collection indexes and archived source
files; they are not generated by the language model. The source URL points to
the original web page or PDF when it is available, whereas the internal SIAGE
action sheets are identified as such and do not receive an artificial public
URL. The original title preserves the wording supplied by the source and makes
it possible to distinguish the source document from a normalised or extracted
metadata value.

At the time of analysis, all 1,143 available LLM records contained a
standardised source identifier. A source URL was available for 1,083 external
records; the remaining 60 records correspond to SIAGE internal action sheets.
An original title was available for 1,142 records, with one RFVAA record lacking
a title in its source index. These coverage figures make missing provenance
visible rather than replacing it with generated or inferred information.

The historical collection indexes did not consistently record a reliable
collection timestamp. The \texttt{collection\_date} field is therefore left
empty for these legacy records rather than being reconstructed from file-system
dates, which would be misleading. Future collection scripts should record the
timestamp automatically at download time, using an ISO 8601 date. Together,
these provenance fields support attribution, verification against the original
document, correction of metadata, and responsible reuse of externally collected
materials.
